\section{The combinatorial symmetry group of the Birkhoff polytope}
Let $D\colon G \to \GL(d,\reals)$ be a faithful representation
and let $P(D) = \conv\{ D(g) \mid g\in G \}$
be the corresponding representation polytope.
Then the vertices of $P(D)$ correspond to the elements of $G$.
We may thus view the affine and combinatorial symmetries
as permutations of $G$ itself.

\begin{lemma}\label{l:reppoltrivsyms}
  Let 
  $D\colon G\to \GL(d,\reals)$
  be a faithful representation
  and $P(D)$ the representation polytope.  
  Then the affine symmetry group $\AGL(P(D))$ 
  as permutation group on $G$ contains 
  $\Gamma(G)$ as defined in the last section.
\end{lemma}

\begin{proof}
  The left multiplications $\lambda_g$ are realized
  by left multiplication with 
  $D(g)$, and the right multiplications $\rho_g$
  by right multiplication with $D(g)^{-1}$.
  If $D$ is an orthogonal representation, then the 
  permutation $g\mapsto g^{-1}$ is realized by transposing 
  matrices, sending $D(g) $ to $D(g)^t = D(g^{-1})$.
  The general case (which we will not need)
  can be reduced to the orthogonal 
  case~\cite[Prop.~6.4]{FrieseLadisch16}.
\end{proof}



Now let $P\colon G= S_n \to \GL(n,\reals)$ be the 
standard permutation representation of 
the symmetric group~$S_n$, and let
\[ B_n := \conv\{ P(\sigma) \mid \sigma\in S_n\}
\]
be the Birkhoff polytope.
Theorem~\ref{main:birkhoff_symgrp} claims that 
$\Gamma(S_n)$ is the combinatorial symmetry group of
$B_n$.
(The second claim of Theorem~\ref{main:birkhoff_symgrp} is that
these symmetries come from isometries of the matrix space.
This is then clear, since
the symmetries in $\Gamma(S_n)$ even act by permuting coordinates of
the matrices.)

\begin{proof}[Proof of Theorem~\ref{main:birkhoff_symgrp}]
  Recall that the Birkhoff polytope consists of the doubly
  stochastic matrices~\cite[Corollary~1.4.14]{LovaszPlummer86}.
  In particular, for each index pair $(i,j)$,
  the equality
  $a_{ij}=0$ describes a facet of the Birkhoff polytope.
  Thus its facets, as subsets of $S_n$, are given by
  the $n^2$ subsets
  \[ F_{ij} = \{\pi\in S_n \mid \pi(i)\neq j \},
     \quad \text{$i$, $j=1$, $\dotsc $, $n$}.
  \]
  It will be more convenient to work with the complements
  \[ A_{ij}= S_n \setminus F_{ij}
           = \{ \pi \in S_n \mid \pi(i)=j \}
  \] 
  of the facets.
  For $\sigma$, $\tau\in S_n$, 
  we have 
  $\sigma A_{ij} \tau^{-1} = A_{\tau i, \sigma j} $.
  We also have 
  $ A_{ij}^{-1} := \{ \pi^{-1} 
                      \mid \pi \in A_{ij}
                   \}
                 = A_{ji}$.
  Moreover, for $i$, $j$, $k$ and $l$ in $\{1,\dotsc,n\}$
  we have
  \[ \card{A_{ij}\cap A_{kl}}
      =
      \begin{cases}
        (n-1)!, & \text{if } i=k, j=l, \\
        0       & \text{if } i=k, j\neq l,\\
        0       & \text{if } i\neq k, j=l, \\
        (n-2)! \quad & \text{otherwise}.
      \end{cases}
  \]  
  Any combinatorial symmetry $\alpha$ permutes the facets and thus
  the sets $A_{ij}$, and preserves cardinalities of their
  intersections.
  
  Let $\alpha\colon S_n \to S_n$ be an arbitrary
  combinatorial symmetry of the Birkhoff polytope.
  We have to show that $\alpha\in \Gamma(S_n)$,
  the group containing the maps
  $\pi \mapsto \sigma \pi^{\pm 1} \tau^{-1}$.
  After replacing $\alpha$ by $\gamma \circ \alpha$
  for some $\gamma\in \Gamma(S_n)$ of the form
  $\gamma(\pi) = \sigma \pi \tau^{-1}$,
  we may  assume that $\alpha(A_{11})= A_{11}$.
  Then $\card{\alpha(A_{12})\cap A_{11}} = \card{A_{12}\cap A_{11}}=0$,
  and thus either $\alpha(A_{12})= A_{1j}$ for some $j\neq 1$
  or $\alpha(A_{12})=A_{j1}$ for some $j\neq 1$.
  If the latter is the case, we compose $\alpha$
  with the map $\pi\mapsto \pi^{-1}$,
  so we may assume that $\alpha(A_{12})=A_{1j}$.
  
  Multiplying $A_{1j}$ from the left with the transposition
  $(2,j)$ yields the set $A_{12}$,
  and so we can assume that $\alpha(A_{12})= A_{12}$.
  
  Now for $j\geq 3$, the set $\alpha(A_{1j})$ has empty intersection
  with $A_{11}$ and $A_{12}$ and thus
  $\alpha(A_{1j}) \in \{ A_{1k}\mid k\geq 3\}$.
  Thus $\alpha$ induces a permutation 
  $\sigma$ of $\{3,\dotsc, n\}$ defined by
  $\alpha(A_{1j})= A_{1,\sigma j}$.
  Thus $\sigma^{-1} \alpha(A_{1j})= A_{1j}$,
  and we may assume that $\alpha(A_{1j})=A_{1j}$
  for all $j$.
  Similarly, we can assume that
  $\alpha(A_{j1})= A_{j1}$ for all $j$.
  
  Thus, after composing $\alpha$ with suitable elements
  from $\Gamma(S_n)$, 
  we may assume that $\alpha$ leaves each of the sets
  $A_{1j}$ and $A_{j1}$ invariant.
  For $k\geq 2$, $l\geq 2$ we have that
  $A_{kl}$ is the unique set $S$ among the sets
  $A_{ij}$ (with $i\geq 2$, $j\geq 2$)
  such that $S\cap A_{k1}=\emptyset = S\cap A_{1l}$.
  It follows that
  $\alpha(A_{kl})=A_{kl}$ for all $k$, $l$.
  Thus $\alpha$ is the identity.
  It follows that the original $\alpha$ was already in 
  $\Gamma(S_n)$. 
\end{proof}

